%! TEX root = main.tex

생성형 AI 서비스의 등장 이후로 전세계의 학계와 산업계가 모델을 학습할 대량의 데이터를 요구하기 시작하였다. 그렇게 다양한 분야에서 전에 없던 수준의 고품질 데이터에 대한 수요가 폭증하자, 데이터는 하나의 전략 자산과도 같은 위상을 지니게 되었다. 결과적으로 기존에 데이터를 독점하고 있던 각 분야의 플랫폼 기업들이 막강한 영향력을 떨칠 수 있게 되었지만 정작 데이터의 생성에 기여한 개인들은 아무런 보상도 받지 못한 채, 기업들의 신뢰할 수 없는 인프라에 의존하게 되는 상황이 지속되었다. Web 3.0은 이러한 데이터 자산 시대에 소유권 구조를 바로잡기 위한 비전으로 각광받기 시작했고, 많은 공학자들이 이를 위한 기술 개발에 몰두하였다. 특히나 많은 연구들이 데이터 원작자의 소유권을 보호하면서 자유롭고 안전하게 데이터를 거래할 수 있는 시스템을 구현하는 것을 목표로 이루어졌다. 본 연구는 그러한 시도의 일환으로써, 영지식 증명 기술을 이용해 데이터의 원문이 아닌, '증거'를 검색하는 시스템을 제안한다. 그 중에서도 수치로 표현되는 의료 데이터를 범위를 기반으로 검색할 수 있는 시스템을 제안함으로써, 환자 개개인이 자신의 의료 데이터가 유출될 걱정 없이 플랫폼에 증거를 업로드, 데이터 수요자는 이를 검색하여 필요한 범위에 적합한 데이터의 존재를 확인할 수 있게 한다. 이 시스템은 zk-SNARK 기반의 영지식 증명 데이터만을 중앙집중 서버에 저장하여 개인정보 유출의 위험을 제거하며, 동시에 수요자가 원하는 범위의 데이터를 효율적으로 검색할 수 있는 시스템을 제안한다.